\subsection{BPatch_memoryAccess}

\label{bkm:Ref160278502}\label{bkm:Ref160277443}\index{BPatch_memoryAccess : : : : : : :!}{

This class encapsulates a memory access abstraction. It contains information that describes the memory access type:

read, write, read/write, or prefetch. It also contains information that allows the effective address and the number of

bytes transferred to be determined.}





\bigskip



{

bool isALoad\index{isALoad_NP: : : : : : :!}()}



{

Return true if the memory access is a load (memory is read into a register).}



{

bool isAStore\index{isAStore_NP: : : : : : :!}()}



{

Return true if the memory access is write. Some machine instructions may both load and store.}



{

bool isAPrefetch_NP\index{isAPrefetch_NP: : : : : : :!}()}



{

Return true if memory access is a prefetch (i.e, it has no observable effect on user registers). It this returns true,

the instruction is considered neither load nor store. Prefetches are detected only on IA32.}



{

short prefetchType_NP\index{prefetchType_NP: : : : : : :!}()}



{

If the memory access is a prefetch, this method returns a platform specific prefetch type.}



{

BPatch_addrSpec_NP getStartAddr_NP\index{getStartAddr_NP: : : : : : :!}()}



{

Return an address specification that allows the effective address of a memory reference to be computed. For example,

on the x86 platform a memory access instruction operand may contain a base register, an index register, a scaling

value, and a constant base. The BPatch_addrSpec_NP describes each of these values.}



{

BPatch_countSpec_NP getByteCount_NP\index{getByteCount_NP: : : : : : :!}()}



{

Return a specification that describes the number of bytes transferred by the memory access.}
